% ED26-BM-RAT2021-6-4: Bonnefon's exact 2021 presentation.
\chapter{Utility Conditionals}
\label{sec:p19-rationality-6-4}

A conditional can concern something an agent wants or wants to avoid.
Interpreting it then involves more than assigning truth values to its two
clauses. Bonnefon's treatment distinguishes using such a sentence to persuade,
using it as a premise of a standard inference, and using it to predict an
agent's behavior or infer the agent's beliefs and preferences. Its central
claim is that practical assumptions about agents can influence both the
credibility of a premise and conclusions that the conditional alone does not
entail. These are different contributions of utility, not a replacement truth
table for material implication \cite[Secs.~1--4, pp.~419--422]{Bonnefon2021Utility}.

\section{What is valued, by whom, and under whose control?}
\label{sec:p19-r64-utility-meaning}

\begin{hbdefinition}{def:p19-r64-utility-conditional}{Utility conditionals and their interpretation}
A \emph{utility conditional} has the form ``if \(p\), then \(q\)'',
where the realization of at least one clause affects an agent's goals:
an action, outcome or state is desirable or undesirable to that agent.
The value is agent-relative. It is not the truth value of the clause,
the frequency of an inference response, or necessarily a monetary amount.
The source's economic use of utility concerns valued consequences;
Bonnefon's underlying formulation describes it as goal achievement.
A \emph{disinterested} use brackets these goals and asks what follows from
the supplied information. A valued conditional can be used in that way;
its subject matter alone does not require an action recommendation
\cite[pp.~888--891]{Bonnefon2009UtilityTheory}.
\end{hbdefinition}

The chapter's opening examples connect washing a car with permission to
borrow it, installing an application update with a phone failure, testifying
against a speaker with an accident, allowing an activity with the speaker's
being fired, Peyton's relocation to Paris with Alex's unhappiness, and studying
instead of partying with good grades. In each, identify the affected agent
before inferring a preference. A benefit to the listener is not automatically
a benefit to the speaker; Alex's unhappiness is not automatically Peyton's
cost. Some sentences perform familiar speech acts, some perform less easily
named acts, and some merely describe a valued situation.

\begin{hbconvention}{def:p19-r64-control-valence}{The chapter's two-axis speech-act comparison}
For a consequent valued by the listener, the prototype classification is:
\[
\begin{array}{c|cc}
 &\text{positive for listener}&\text{negative for listener}\\\hline
 \text{speaker controls consequence}&\text{promise}&\text{threat}\\
 \text{speaker does not control consequence}&\text{tip}&\text{warning}
\end{array}
\]
This reconstructs Figure~6.4.1's control/valence comparison, not four
necessary-and-sufficient semantic definitions. The relevant control and
valence are interpreted, not read off from the word ``if''.
\end{hbconvention}

An apparent prediction of an accident after testimony can communicate a
veiled threat although it does not explicitly say that the speaker causes
the accident. Similarly, the speaker may call a threat a promise. The chapter
reports eye-movement work by Wray and colleagues (2016), Haigh and colleagues
(2011, 2014), Stewart and colleagues (2013), and Wood and colleagues (2016)
as showing graded boundaries, especially around threats. The explicit
control distinction does not invariably determine the pragmatic category.
These are attributed review-level findings: this chapter supplies no common
sample, pooled processing-time estimate or universal classifier for that
set of reading experiments \cite[pp.~419--420]{Bonnefon2021Utility}.

\section{Three reasoning tasks, not one correctness criterion}
\label{sec:p19-r64-tasks}

The persuasive task points to good or bad consequences in support of a
course of action or policy. Its practical question is whether the reasons
support the proposed choice, including the probabilities, alternatives and
values assumed. Bonnefon connects it with arguments from consequences and
their Bayesian/argumentation literature, but explicitly does not develop that
research in this chapter. Persuasiveness is not thereby either a proof of
the consequence's truth or a proof that the choice maximizes utility.

The standard-inference task uses the conditional with a categorical premise.
In the Peyton/Alex example, a move implying unhappiness and Alex's being happy
invite a modus-tollens conclusion that Peyton did not move. This formalizes
the source's example only if happiness is stipulated to exclude the specified
unhappiness, the same time and agents are involved, and the conditional is
accepted with the relevant interpretation. The Lucy example instead uses
modus ponens: failure to eat salad, together with a five-year grounding
threat, yields that grounding under a fixed material reading. Doubting the
threat does not construct a countermodel to modus ponens.

The nonstandard-inference task starts with the sentence connecting Alex's
leaving early to her dismissal, then predicts that she will not leave.
Neither \(p\) nor \(\neg p\) follows from \(p\to q\) alone:
the valuation \(p=q=1\) refutes the proposed \(\neg p\) conclusion,
and \(p=q=0\) refutes the proposed \(p\) conclusion, while both make
the conditional true. The behavioral prediction adds assumptions about what
Alex values, knows and can choose. Calling it \emph{nonstandard} or
\emph{paralogical} distinguishes that route from deductive consequence;
it does not establish that a sensible defeasible prediction is irrational
\cite[p.~420]{Bonnefon2021Utility}.

\Cref{thm:p19-conditional-benchmark} supplies the fixed-meaning deductive
benchmark. \Cref{sec:p19-rationality-6-2} supplies an explicitly suppositional
probability assessment. A probability of \(q\) supposing \(p\), a judgment
that an inference is necessary, and a prediction that an agent chooses
\(p\) have different targets. Good performance on one does not answer
the others.

\section{Credibility can carry practical assumptions into standard inference}
\label{sec:p19-r64-standard}

Bonnefon contrasts a parent offering a favorite dessert after a room is tidied
with the advice that hard work will obtain a good grade from a teacher.
The parent's control over the first outcome can support greater confidence
in fulfillment; the adviser cannot guarantee the second. The chapter reports
more MP endorsement from promises than tips, citing Couto--Quelhas--Byrne
(2017) and Evans and colleagues (2008), and gives the analogous threat/warning
comparison. Control provides one reason for assessing credibility differently,
not a proof that every promise is more credible than every tip. Ability,
sincerity, competing commitments and knowledge of the situation can reverse
a particular comparison \cite[pp.~420--421]{Bonnefon2021Utility}.

Credibility can also depend on the enforcer's incentives. Grounding a child
for an evening after refusing salad differs from grounding the child for five
years. The latter consumes the parent's options for later sanctions and
requires sustained enforcement effort. Assuming some practical rationality
in the parent can therefore reduce confidence that the sanction will occur.
The inference's formal shape has not changed. What has changed is acceptance
or confidence in its main premise, drawing on a model of another agent's
practical behavior. The chapter cites Amgoud--Bonnefon--Prade (2005),
L\'{o}pez-Rousseau and colleagues (2011), and Verbrugge and colleagues (2004)
for this credibility question, not a new numerical enforcement experiment.

Exceptions contribute another distinction. The reviewed legal tasks compare
unlicensed music downloading with killing, and punishment for evading an
event's admission charge with punishment for maltreating a minor in one's
care. The chapter reports easier generation of exceptions for the downloading
example and correspondingly less MP endorsement in predictive sentences
using ``will''. With ``should'', it reports that exculpatory circumstances
affect the less morally charged example more than the maltreatment example,
and interprets this difference through moral outrage. The selected
Gazzo Casta\~neda--Knauff (2016a,b) and Gazzo Casta\~neda--Richter--Knauff
(2016) references are review-level evidence here: no original recruitment,
exclusion count, endorsement denominator or confidence interval is provided
by the 2021 chapter. The claimed emotion mechanism is its interpretation,
not identified merely by observing a difference in endorsement.

These examples ask about interpretation and judgment, not about the actual
content of current law. ``Will be punished'' predicts an outcome; ``should
be punished'' expresses a normative assessment. Replacing one modal with the
other changes the question. A participant's refusal to consider an exception
does not establish that the exception is legally or morally irrelevant,
and the validity of MP does not justify a sanction. Formal validity,
premise credibility, observed judgment and normative warrant remain separate
\cite[p.~421]{Bonnefon2021Utility}.

\section{Nonstandard predictions and their additional premises}
\label{sec:p19-r64-nonstandard}

The application/award example predicts that Alex will apply, provided the
award matters enough to outweigh the effort, she knows that applying obtains
it, she is able to apply, and her choice responds to that advantage. Increasing
the award's value, lowering its probability conditional on applying, or
making a cheaper route salient need not have the same effect. The 2021
chapter proposes stronger action predictions with greater desirable value,
weaker predictions with lower conditional probability, and weaker predictions
when a less costly alternative achieves the goal. It also cites
Haigh--Bonnefon (2015a,b) for rapid reading-time sensitivity. These are
source-reported directions, not a population response function or evidence
that participants explicitly calculate expected utilities.

\begin{hbdefinition}{def:p19-r64-binary-choice}{A finite decision bridge with its assumptions visible}
As an explicit mathematical illustration, an agent can choose action
\(a_1\) or alternative \(a_0\). Their probabilities of a valued outcome
are \(r_1,r_0\in[0,1]\) under the consequences of those choices.
Relative to not receiving the outcome, receiving it adds utility \(v\in\mathbb R\).
The action costs are \(c_1,c_0\in\mathbb R\), assessed on the same additive scale.
The two expected utilities, up to a common constant, are
\[
 U_1=r_1v-c_1,\qquad U_0=r_0v-c_0.
\]
Assume these consequences include everything relevant, the probabilities
and utility scale are the agent's, and choice maximizes this expectation
over these two available acts. This is a stipulated application of the
P17 framework, not a numerical model supplied or fitted by the chapter.
\end{hbdefinition}

\begin{hblemma}{thm:p19-r64-choice-gap}{proved=annex:p19.r64.choice-gap}{What a utility conditional does not determine by itself}
Under \cref{def:p19-r64-binary-choice}, action \(a_1\) is uniquely preferred
exactly when \((r_1-r_0)v>c_1-c_0\), and both acts maximize utility
when equality holds. For fixed probabilities and costs, increasing \(v\)
strengthens the advantage of \(a_1\) only if \(r_1>r_0\); it weakens
that advantage if \(r_1<r_0\), and leaves it unchanged if \(r_1=r_0\).
For fixed \(v>0\), increasing \(r_1\) strengthens its advantage.
Knowledge of \(r_1\) and \(v\) alone does not determine choice without
the alternative and cost assumptions.
\end{hblemma}

For an illustrative award worth \(v=10\), application success
\(r_1=.8\), alternative success \(r_0=.2\), and costs
\(c_1=3,c_0=0\), the values are \(U_1=5\), \(U_0=2\),
so applying wins. At \(r_1=.4\), they become \(1\) and \(2\),
so applying loses. If instead a free alternative has success probability
\(.9\), its value is \(9\), exceeding the original application's \(5\).
These declared examples justify the directions and their boundary without
asserting that a desirable consequent by itself makes its antecedent likely.
The utility scale is not merely the ordinal scale of
\cref{sec:p17-finite-preferences}; averaging it invokes the stronger
domain and assumptions of \cref{sec:p17-finite-expected-utility}.

The underlying 2009 account represents the relevant information in a
\emph{utility grid}. The antecedent row records its actor, utility sign and
affected agent; the consequent row records the corresponding triple:
\[
 \begin{array}{c|ccc}
 &\text{actor}&\text{utility}&\text{target}\\\hline
 \text{if-clause}&x&u&y\\
 \text{then-clause}&x'&u'&y'
 \end{array}
\]
Here \(u,u'\in\{-,0,+\}\) are signs, not cardinal amounts.
\(s,h,e\) name speaker, hearer and another identified agent;
\(\omega\) is a neutral world-agent convention for a nonintentional event.
\(A\) is the set of agents: zero known utility is recorded with target
\(A\). That coding does not establish actual indifference or absence of
unmentioned goals. Multiple affected agents can require multiple subgrids.
The grid records theorists' information, not an asserted literal mental
array \cite[pp.~891--894]{Bonnefon2009UtilityTheory}.

The proposed \emph{folk axioms} are beliefs reasoners are hypothesized to use:
self-interested behavior predicts that agents act for their own benefit and
avoid their own loss; self-interested attitude predicts that an agent favors
another's action that benefits that agent, and opposes one that harms that
agent; limited altruism predicts benefit to others when it does not cost the
actor, and avoidance of harm to others when it does not benefit the actor.
They are neither axioms of logic nor demonstrated universal laws of choice.
The 2009 source presents a nonexhaustive qualitative framework and explicitly
does not supply a complete quantitative Bayesian translation. Its limited-
altruism evidence is identified as indirect. In particular, predicting what
the speaker thinks another should do is not endorsing that recommendation
oneself \cite[pp.~894--898]{Bonnefon2009UtilityTheory}.

\section{Unexpected behavior: infer a missing belief or a different preference?}
\label{sec:p19-r64-revision}

In the loyalty-card example, Brenda could save twenty dollars by using the
card but does not use it. The ordinary expectation depends on the saving's
importance, her knowledge of it, a suitably goal-responsive decision process,
and her ability to use the card. If ability and decision responsiveness remain
assumed, ignorance of the saving and indifference to that amount are rival
explanations. The behavior alone does not determine which premise to revise.
Nor should an observer retain ability by default when the card is unavailable.

Bonnefon--Girotto--Legrenzi's selected original experiment supplies an important
boundary to this account. Campus recruits at Toulouse answered in French.
The methods print 62 women and 16 men, mean age~21, SD~4.7, but print group
sizes 35 for prediction and 37 for revision, totaling~72 rather than~78.
The prediction tests later use \(df=39\) and the revision tests use
\(df=37\). These are inconsistent published sample indicators; none is
silently repaired into a verified denominator. The paper reports random
assignment to the two groups, four brief saving scenarios, and four rotated
amounts, 40, 80, 160 and 320 euros. Its predictions, knowledge assessments
and importance assessments used eleven-point ratings from \(-5\) to \(+5\),
not observed spending behavior \cite[pp.~34--36]{BonnefonGirottoLegrenzi2012Preferences}.

The prediction group expected the money-saving action and attributed both
knowledge and importance. At forty euros, its reported means for these three
ratings were \(2.9,2.6,1.4\); at 320 euros they were \(3.4,2.0,3.2\).
The revision group was instead told that the action was not taken and rated
knowledge and importance. Its corresponding means were \((1.3,-.9)\)
at forty euros and \((-.1,.4)\) at 320 euros. These are points on the
rating scale, not probabilities or percentages of people retracting a belief.
The study did not test the full set of alternative explanations: the folk-
axiom assumption was retained and questions offered knowledge/importance
revision \cite[pp.~35--39]{BonnefonGirottoLegrenzi2012Preferences}.

For forty euros the authors interpret the pattern as retaining knowledge
while discounting importance. For larger amounts, near-zero group means
conceal opposed individual patterns rather than establish universal ignorance.
Across revision participants' scenario-averaged ratings, the reported
knowledge/importance correlation was \(r(38)=-.54,p=.001\).
No participant-level interval is supplied here, and the denominator
inconsistency remains material to exact inferential reconstruction.
Neither the correlation nor the four scenarios establish a stable personality
trait, a uniquely rational revision rule, or longitudinal within-person
belief change: prediction and revision were different groups. The article
itself calls for further psychometric work \cite[pp.~37--40]{BonnefonGirottoLegrenzi2012Preferences}.

\section{Templates constrain the assigned utilities}
\label{sec:p19-r64-templates}

The chapter's three-agent example connects your hurting the speaker to the
speaker's helping a third person. One interpretation supplies two positive
values: the speaker welcomes the supposed hurt, while you welcome the
speaker's assistance to the third person. Another supplies two negative
values: the speaker dislikes the hurt and you dislike that assistance, so
the offer functions as retaliation. No one interpretation follows from the
sentence alone. Fitting a familiar reciprocal pattern can change the
preferences attributed to the agents \cite[p.~422]{Bonnefon2021Utility}.

The Peyton/Jesse example uses unknown verbs for an action and its return:
Peyton tymps Jesse, then Jesse yorbs Peyton. In a same-sign social-contract
interpretation, Jesse's enjoying the first action supports Peyton's enjoying
the return; disliking the first supports disliking the return. This is a
template hypothesis, not logical information contained in the unknown verbs.
Similarly, if Peyton hurts Jesse and Jesse zims Peyton, a retaliatory reading
suggests that Peyton dislikes being zimmed. Supplying that Peyton likes it
can prompt a reinterpretation in which Jesse likes the apparently harmful
action. The semantics of ``hurt'' has not become the proposition that every
target enjoys injury; the question is context-specific attributed value.

\begin{hbdefinition}{def:p19-r64-templates}{Four proposed utility templates}
For nonzero sign \(u\), the Social Contract has rows
\((x,u,y),(y,u,x)\); Justice has rows \((x,u,y),(x,u,x)\).
Both Unpacking and Tautology have the actor/target/sign arrangement
\((x,u,y),(x,u,y)\), but they differ in explanation: Unpacking makes
the consequent the valued result implicit in the first action; Tautology
restates the first action in another description. The grid alone cannot
distinguish them. These are proposed interpretive schemas, not four
Boolean connectives \cite[pp.~351--354]{BonnefonHaighStewart2013Templates}.
\end{hbdefinition}

The original completion survey involved 30 Manchester students, 18 women,
mean age~27.5, SD~8.6, all native English speakers. Each completed thirty
fragments in an individually randomized order, giving 900 completions.
Justice, Social Contract and Unpacking accounted for reported 13\%, 11\%
and 10\% of completions. Tautology accounted for 2\% but was selected
for structural simplicity. These are completion shares in a restricted
fragment survey, not estimates that those shares govern ordinary conversation
or that no other templates exist.

Experiment~1 recruited a different 60 Manchester participants, 54 women,
for course credit, using online survey software. Sixteen sentences with
four-letter nonverbs crossed four template structures and positive/negative
information about a companion nonverb. A seven-point \(-3\) to \(+3\)
rating asked about the target verb's intrinsic utility on other occasions,
to exclude liking an action merely because of this trial's downstream
consequence. Social-contract target means were \(1.6\) (SD~1.1) after a
positive companion and \(-1.5\) (SD~1.3) after a negative one.
The companion-valence effect was \(F(1,59)=230.8,p<.001\), reported
\(\eta_p^2=.80\); valence interacted with template,
\(F(3,57)=4.0,p=.012,\eta_p^2=.17\). Such utility ratings
do not identify a unique reading algorithm or certify a normative inference
\cite[pp.~354--355]{BonnefonHaighStewart2013Templates}.

Experiment~2 used a different 56 Manchester participants, 49 women, again
for course credit online. Sixteen items crossed template, the real verbs
help/hurt, and a control versus conflicting-value condition. The three-way
interaction was \(F(3,53)=13.2,p<.001,\eta_p^2=.43\).
Congruency effects were reported for help in every template; for hurt,
Justice was nonsignificant, \(t(55)=1.34,p=.19\), while the other
three showed effects. Therefore the chapter's illustration should not be
read as a universal reversal for every known verb or template. The authors
also noticed near-neutral control ratings for help/hurt and discussed demand
characteristics. Their later experiments address these concerns and mere
sign matching; they are not added to this bounded numerical record as
independent replications. No original-study exclusion count or new confidence
interval is manufactured \cite[pp.~355--356]{BonnefonHaighStewart2013Templates}.

\section{A consequence model cannot silently become a diagnostic model}
\label{sec:p19-r64-causal-boundary}

The source's action predictions assume that acting makes the valued outcome
more or less likely. A sentence can instead describe a symptom: accepting
a deal could cause ruin, or could indicate that the person is already ruined
and therefore attracted to a desperate offer. A high observational
\(P(q\mid p)\) is not by itself the probability of \(q\) under choosing
\(p\). This qualification is explicit in Bonnefon--Sloman's source cited
by the 2021 chapter \cite[pp.~195--197]{BonnefonSloman2013Causal}.

\begin{hblemma}{thm:p19-r64-diagnostic-failure}{proved=annex:p19.r64.diagnostic-failure}{Perfect diagnosis need not make an action costly}
In a finite structural model let a background indicator \(Z\) be one
with probability \(1/2\), let the observed action \(D=Z\), and let
the undesirable outcome \(Q=Z\). Then \(P(Q=1\mid D=1)=1\),
but setting \(D\) to either zero or one leaves the probability of
\(Q=1\) equal to \(1/2\). If choosing \(D=1\) adds direct benefit
one and \(Q=1\) incurs loss ten, the two intervention-based expected
values are \(-4\) and \(-5\), respectively. The beneficial action
therefore wins despite perfect diagnostic association with the bad outcome.
\end{hblemma}

The construction is a declared failure case, not a fitted model of the
deal experiment. It uses the equation-replacement semantics of
\cref{def:p16-causal-intervention}; it does not infer that people use
structural equations or that a threatening utterance merely reports a symptom.

Bonnefon--Sloman's Experiment~4 recruited 53 Mechanical Turk participants,
19 men and 34 women, mean age~33, SD~11. Thirty-two scenarios in randomized
order crossed causal/diagnostic structure with clear/unclear utility within
participants. They asked whether a man thought a woman's action should
be taken; they did not ask whether participants endorsed the action themselves.
A causal gloss with ``will'' contrasted with a diagnostic gloss with
``must''. Unclear utility named changed prospects without their direction.
The eleven-point \(-5\) to \(+5\) ratings were sign-reversed on
negative-consequence items so that higher index values marked stronger
utility-consistent attitude attribution.

For clear utility the reported causal/diagnostic means were \(3.7,2.8\),
\(t(52)=5.1,p<.001\), difference 95\% CI \([.5,1.2]\),
reported \(d=.53\). For unclear utility they were \(-.2,-.1\),
\(t(52)=-.8,p=.42\), CI \([-.3,.1]\), reported \(d=.12\).
The interaction was \(F(1,52)=26,p<.001\). The source labels its
mean-error quantities MSE; no different precision statistic is inferred
from that label. The nonsignificant unclear-utility contrast is not an
equivalence proof. Note~7 acknowledges diagnostic utility effects and
possible normative/necessity readings of ``must''. The result therefore
supports a qualified contrast under the manipulated language package,
not zero diagnostic influence or an isolated causal-representation mechanism
\cite[pp.~201--203,206, note~7]{BonnefonSloman2013Causal}.

\section{Scope, other-regarding preferences, and consultation}
\label{sec:p19-r64-consultation}

Bonnefon concludes that utility conditionals provide an interface between
reasoning, decision research, theory of mind, Bayesian accounts and moral
judgment. Its future direction is other-regarding preferences: how actors
weight others' outcomes, and how observers think they do. The chapter's
claim of an undeveloped theory is the research gap identified in this
presentation, not a theorem that no economic model of such preferences
exists. Its references to Cooper--Kagel (2015), De Vito--Bonnefon (2014)
and Jara-Ettinger and colleagues (2016) identify relevant neighboring work,
not a complete weighted-social-utility model supplied here
\cite[p.~422]{Bonnefon2021Utility}.

\begin{hbmethod}{def:p19-r64-consultation}{Using consequences without confusing the task}
First specify whether the requested conclusion is a fact, confidence,
behavioral prediction, attributed preference, attributed recommendation,
or recommendation endorsed by the observer. Identify each actor, target,
value and perceived control. For a standard inference fix the conditional's
meaning and premise standing. For an action prediction state the agent's
knowledge, ability, alternatives, costs, outcome probabilities and decision
assumption. Check whether the stated link is consequential or diagnostic.
If an expected action is not observed, compare missing belief, changed value,
inability and a different decision assumption rather than select one by fiat.
Keep any template-based reinterpretation visible and defeasible.
\end{hbmethod}

A success case is the declared two-act award calculation, or explaining
which stated assumptions make a reciprocal unknown-verb interpretation
plausible. Failure cases are inferring \(\neg p\) from the conditional
alone, inferring an action from a benefit without its alternative, treating
perfect diagnosis as an effect of acting, converting an attributed ``should''
into an observer's normative endorsement, or reading a near-zero group
mean as identical beliefs in every participant.

\Cref{sec:p19-conditional-reasoning-suppression,
sec:p19-conditional-interpretation-experiments} distinguishes premise
interpretation from unchanged formal validity under extra information.
\Cref{sec:p19-rationality-6-3} treats possibility representations and
counterfactual tasks; templates for attributed value are not those possibility
rows. The suppositional assessment of \cref{sec:p19-rationality-6-2}
does not by itself supply an agent's cost or preferences, and the P17
representation theorem does not by itself establish observed compliance.
The present chapter's contribution is the specified route by which those
practical assumptions enter interpretation and prediction, together with
its task-specific evidence and limitations.
